\chapter{Conclusion and Perspectives}

\section{Conclusion}
The importation market in Algeria is vast and full of potential, yet it has long been hindered by informal practices, a lack of transparency, and significant security risks for clients. Through this graduation project, we successfully addressed these challenges by designing, developing, and deploying \textbf{Sila DZ}, an innovative and intelligent platform that connects clients and importers in a secure ecosystem.

By leveraging a robust Microservices architecture—powered by Spring Boot, NestJS, and Kafka—we created a highly scalable backend capable of handling real-time negotiations and complex order tracking. On the frontend, the combination of React.js for the web and Flutter for mobile ensured a premium, accessible, and RTL-compliant user experience. We successfully translated a theoretical business model, which includes a dynamic Points System and rigorous identity verification, into a tangible software product. Sila DZ not only digitizes the importation workflow but fundamentally builds trust, democratizing micro-importation for the Algerian population.

\section{Perspectives}
While the current version of Sila DZ serves as a powerful Minimum Viable Product (MVP) and a solid foundation for the startup, technology and business needs are always evolving. We have identified several key perspectives for future development:
\begin{itemize}
    \item \textbf{Advanced Payment Gateway Integration:} Integrating local banking APIs (CIB/Edahabia) and international payment solutions to fully automate financial transactions and currency conversions directly within the platform.
    \item \textbf{Logistics \& GPS Tracking:} Partnering with international freight forwarders to provide live GPS tracking of shipping containers, allowing clients to track their goods on a map in real-time.
    \item \textbf{Artificial Intelligence (AI):} Implementing Machine Learning algorithms to provide smart recommendations, automatically matching clients with the most suitable importers based on historical data, ratings, and pricing trends.
    \item \textbf{Bidding System for Large Enterprises:} Expanding the business model to allow large companies to post tenders (Appels d'offres), where registered importers can submit competitive bids in a structured auction format.
\end{itemize}

Ultimately, this project represents more than an academic thesis; it is the launchpad for a real-world startup poised to modernize the Algerian supply chain and connect local entrepreneurs with global opportunities.
